\subsection{The exact normalization of the Tunnell coefficient}
\label{sec:tunnell-normalization}
Write
\[
 \frac{\eta(4\tau)^5\eta(16\tau)}{\eta(2\tau)^2\eta(8\tau)}
   =\sum_{n\geq1}a(n)q^n,
 \qquad
 \varphi_{8|4}(\tau,z)=\sum_{n,r}C_{\varphi}(r^2-4n)q^ny^r,
\]
where $\varphi_{8|4}=(y^{-1}+y)q+O(q^2)$ is the normalization in
Cheng--Duncan--Mertens \cite[equation~(4.2.9)]{CDM2025}.
The exact comparison in this normalization is
\begin{equation}\label{eq:tunnell-factor-two}
 C_{\varphi}(-n)=\frac12a(n)\qquad(n\equiv3\pmod8).
\end{equation}
In particular $C_{\varphi}(-3)=1$ whereas $a(3)=2$.
The factor two does not affect a nonvanishing criterion, but does affect
an identity of coefficients and any explicit constant multiplying their
squares.

\begin{proof}
First we identify the Jacobi form itself:
\begin{equation}\label{eq:phi84-explicit}
 \varphi_{8|4}(\tau,z)
   =\eta(2\tau)\eta(8\tau)^2\theta_{1,1}(\tau,z).
\end{equation}
Here is a modular verification, in addition to the leading-term check.
The Jacobi triple product gives
$\theta^0_{1,1}=2\eta(4\tau)^2/\eta(2\tau)$, so the right side is
\[
 2\eta(4\tau)^2\eta(8\tau)^2
    \frac{\theta_{1,1}(\tau,z)}{\theta^0_{1,1}(\tau)}.
\]
The eta criterion puts $\eta(4\tau)^2\eta(8\tau)^2$ in
$S_2(\Gamma_0(32))$ with trivial character: its two congruence sums are
$2\cdot4+2\cdot8=24$ and $2\cdot(32/4)+2\cdot(32/8)=24$,
and $4^2 8^2$ is a square.  The theta transformation law shows that
$\theta_{1,1}$ transforms by a scalar on $\Gamma_0(4)$, with the same
scalar at $z=0$; their ratio has weight zero and index one there.
It is holomorphic in the upper half-plane, since $\theta^0_{1,1}$ has
no zeros there.  Cusp holomorphy can be checked without dividing:
the positive eta product $\eta(2\tau)\eta(8\tau)^2$ has positive order
at every cusp, and the transformed index-one theta vector has only
nonnegative theta exponents.  Their product is therefore a Jacobi cusp
form of weight two and index one for $\Gamma_0(32)$.
Its first term is $(y^{-1}+y)q$.  The uniqueness specified in
Cheng--Duncan--Mertens \cite[equation~(4.2.9)]{CDM2025} proves \eqref{eq:phi84-explicit}.

Put $\vartheta(\tau)=\sum_{r\in\mathbb Z}q^{r^2}$.
Another triple-product identity gives
\[
 \frac{\eta(4\tau)^5\eta(16\tau)}{\eta(2\tau)^2\eta(8\tau)}
   =\vartheta(2\tau)\eta(8\tau)\eta(16\tau).
\]
The eta product on the right has exponents congruent to one modulo eight.
The terms with exponent congruent to three modulo eight come precisely
from the odd integers in $\vartheta(2\tau)$.  As
\[
 \sum_{r\ \mathrm{odd}}q^{2r^2}
    =2\frac{\eta(32\tau)^2}{\eta(16\tau)},
\]
we obtain the exact formal-series identity
\begin{equation}\label{eq:tunnell-sieve}
 \sum_{n\equiv3\ (8)}a(n)q^n=2\eta(8\tau)\eta(32\tau)^2.
\end{equation}
On the other hand the odd theta coefficient of
\eqref{eq:phi84-explicit}, with $\tau$ replaced by $4\tau$, is
\[
 \sum_{n\equiv3\ (8)} C_{\varphi}(-n)q^n
   =\eta(8\tau)\eta(32\tau)^2.
\]
Comparison with \eqref{eq:tunnell-sieve} proves
\eqref{eq:tunnell-factor-two} for every such $n$, without a finite
coefficient extrapolation.
\end{proof}

The coefficient $a(n)$ here is exactly the one in Tunnell's original
definition $g\vartheta_2$, where
$g=q\prod_{r\geq1}(1-q^{8r})(1-q^{16r})$
\cite[p.~323]{Tunnell1983}.  Thus the correction is a factor between the
two normalizations, not a change to Tunnell's definition.
